\section{Significance and Innovation}

\subsection{Disease Burden and Unmet Need}

Early-onset Parkinson disease (EOPD, onset $\leq$50 years) represents approximately 10--15\% of all PD diagnoses yet accounts for a disproportionate share of years of life lost and disability-adjusted life years (DALYs). A 2024 analysis of Canadian administrative health data (\textit{ICES Cohort}, $n = 4{,}812$) confirmed that EOPD patients experience diagnostic delays averaging 3.4 years, accelerated motor progression, and a 2.8-fold higher rate of treatment-resistant dyskinesia compared with age-matched late-onset PD.

Current pharmacological strategies --- primarily levodopa/carbidopa and dopamine agonists --- address symptoms but do not modify the underlying neurodegenerative process. No CIHR-funded multisite trial has yet evaluated neuroinflammation as a therapeutic target specifically in EOPD. This proposal addresses that gap directly.

\subsection{Scientific Innovation}

\begin{tcolorbox}[cihrbox, title={\faLightbulb\enspace Key Innovations}]
  \begin{itemize}[leftmargin=14pt, itemsep=2pt, topsep=2pt]
    \item \textbf{First EOPD-specific microglia atlas}: Leveraging 10x Visium spatial transcriptomics and snRNA-seq on 60 post-mortem substantia nigra specimens stratified by age-of-onset.
    \item \textbf{Novel microglial subtype --- ``EO-MG1''}: Preliminary data (see Fig.\ 1) reveal a disease-associated microglial cluster expressing \textit{GPNMB$^{\text{hi}}$}/\textit{TREM2$^{\text{lo}}$} uniquely enriched in EOPD nigra ($p < 0.001$ vs.\ LOPD and controls).
    \item \textbf{Translatable biomarker pipeline}: Plasma GPNMB and soluble TREM2 quantified across a 320-participant longitudinal cohort in partnership with Apex Health Sciences Centre.
    \item \textbf{Dual therapeutic testing}: CSF1R antagonist (Compound NX-7714, Nexa Biosciences, preclinical licence) and IL-6 trans-signalling blockade tested in humanised microglia-neuron co-culture and AAV-$\alpha$-synuclein rat model.
  \end{itemize}
\end{tcolorbox}

%% ─── SECTION 3: OBJECTIVES & HYPOTHESES ──────────────────────────────────────
